% !PRINTED 9
\weberpage{9}

%<<CONT
\weq{5}{\dfrac{(x_1x_2)(x_3x_4)}{(x_1x_3)(x_2x_4)} = \dfrac{(x_1'x_2')(x_3'x_4')}{(x_1'x_3')(x_2'x_4')}.}

If the function $f(x)$ is given in the elliptic differential

\[
du = \dfrac{dx}{\sqrt{f(x)}}
\]

then the $x_1$, $x_2$, $x_3$, $x_4$ are thereby given as well, and by the linear transformation (3) one can bring the differential into another form in which three of the quantities $x_1'$, $x_2'$, $x_3'$, $x_4'$ have arbitrarily given values. % [sic] German source prints singular "Größe" after "three of the"
One obtains the \gesperrt{normal form} under the following assumptions:

\[
x' = z, \quad x_1' = 0, \quad x_2' = 1, \quad x_3' = \dfrac{1}{\varkappa^2}, \quad x_4' = \infty
\]

The quantity $\varkappa$ is called by \gesperrt{Legendre} the \gesperrt{modulus of the elliptic differential}, and $\sqrt{1-\varkappa^2} = \varkappa'$ the \gesperrt{complement of the modulus}. If, then, one assigns the points to one another in the following way:

\[
\begin{array}{ccccc}
z, & 0, & 1, & \dfrac{1}{\varkappa^2}, & \infty, \\[6pt]
x, & x_1, & x_2, & x_3, & x_4,
\end{array}
\]

then from the equality of the cross ratios the relations follow easily:

\weq{6}{\begin{gathered}
z = \dfrac{(xx_1)(x_2x_4)}{(xx_4)(x_2x_1)}, \\
1 - z = \dfrac{(xx_2)(x_1x_4)}{(xx_4)(x_1x_2)}, \\
1 - \varkappa^2 z = \dfrac{(xx_3)(x_1x_4)}{(xx_4)(x_1x_3)},
\end{gathered}}

\weq{7}{dz = \dfrac{(x_4x_1)(x_4x_2)}{(x_2x_1)}\,\dfrac{dx}{(xx_4)^2},}

\weq{8}{\varkappa^2 = \dfrac{(x_3x_4)(x_1x_2)}{(x_1x_3)(x_2x_4)}, \quad \varkappa'^2 = \dfrac{(x_2x_3)(x_1x_4)}{(x_2x_4)(x_1x_3)}}

and from (3):

\weq{9}{\dfrac{\sqrt{(x_3x_1)(x_4x_2)}\,dx}{\sqrt{-(x_1x)(x_2x)(x_3x)(x_4x)}} = \dfrac{dz}{\sqrt{z(1-z)(1-\varkappa^2z)}}.}

According to § 2, $\varkappa^2$ becomes a positive proper fraction if the $x_1$, $x_2$, $x_3$, $x_4$, are real and follow one another cyclically in order of magnitude. This therefore gives eight different such transformations,
%CONT>>
